\chapter{Efectos biológicos de la radiación}

El núcleo es la parte más sensible a la radiación. A este nivel, la radiación gamma colisiona con los enlaces de la molécula de ADN, alterando el gen en el cual se encuentra la molécula de ADN. Este gen alterado puede llegar a producir un efecto en los cromosomas que, al dividirse, produce células alteradas con funciones diferentes a las de la célula madre. Esto se conoce como efecto genético de la radiación ionizante.

En general, podemos decir que la interacción de la radiación ionizante con las células puede producir en ellas:
\begin{itemize}
  \item aumento o disminución de volumen;
  \item su muerte;
  \item un estado latente;
  \item mutaciones en las células hijas.
\end{itemize}

La radiosensibilidad es la respuesta que las células ofrecen a la acción de la radiación, la cual es directamente proporcional a la capacidad reproductiva de las células y varía inversamente con el grado de diferenciación de las mismas, así como con su estado metabólico, su estado de división y su estado de nutrición.

Por lo tanto, las células que más rápidamente se dividen son las más radiosensibles (hematopoyéticas, piel, cristalino, epitelio, gastrointestinales, gónadas, etc.).

En función de los dos tipos de células que se encuentran en los organismos, los efectos biológicos de la radiación se clasifican en dos tipos: somáticos y hereditarios.

\section{Efectos somáticos y hereditarios}

Los efectos somáticos son aquellos que se manifiestan en el individuo que se ha expuesto a la radiación y, como su nombre lo indica, se deben al daño recibido en sus células somáticas o diploides; además, este daño queda limitado solamente al individuo. Algunos ejemplos de este tipo de alteraciones son la esterilidad, la disminución de células en la médula ósea y la inducción al cáncer.

Los efectos hereditarios o genéticos son el resultado del daño recibido en las células gaméticas y sus efectos se presentarán en la descendencia de las células irradiadas. Un ejemplo serían los procesos mutagénicos en el ADN, los cuales serían heredados a sus descendientes.

\section{Efectos determinísticos}

Actualmente, los efectos biológicos de la radiación, de acuerdo con su probabilidad de incidencia, se clasifican en dos tipos: efectos determinísticos (no estocásticos) y efectos estocásticos.

Los efectos determinísticos se definen como aquellos que se producen a partir de una dosis umbral y aumentan en severidad dependiendo de la dosis recibida. El efecto se agudiza para dosis altas recibidas en un tiempo corto; estos efectos se derivan de la muerte celular.

La dosis umbral para algunos efectos determinísticos se presenta a continuación, cuando la irradiación es a cuerpo total.

\subsection{Exposición aguda a cuerpo total}

\textbf{Dosis absorbida (rads) y efecto:}
\begin{description}
  \item[De 10] Ruptura cromosómica difícil de detectar en células sanguíneas. Interferencia con la organogénesis de embriones.
  \item[De 25] Cambios sanguíneos.
  \item[De 50] Probable retención momentánea de la espermatogénesis.
  \item[De 100] Síndrome de radiación probable.
  \item[De 200] Síndrome grave de muerte.
  \item[De 400] 50\% de probabilidades de muerte.
  \item[De 600] 100\% de probabilidades de muerte.
\end{description}

Los síntomas del síndrome de radiación son: náuseas, vómitos, agotamiento, disminución de glóbulos blancos, período de recuperación aparente, vómito y diarrea con sangre, disminución de glóbulos blancos y plaquetas, anemia y, finalmente, intoxicación general.

Se llama exposición aguda cuando la dosis se recibe en un tiempo corto, por lo cual el cuerpo tiene menor oportunidad de reparación, y exposición crónica cuando la dosis se recibe en un tiempo largo.

El tiempo que tardan en presentarse los síntomas depende de la magnitud de la dosis y del intervalo de tiempo en que se imparta esta (dosis aguda o dosis crónica); es decir, cuando la dosis se imparte en un tiempo mayor, el efecto es menor, como si el organismo recibiera una dosis menor.

En el caso de que la irradiación no sea en todo el cuerpo sino localizada en una parte de él, el daño depende de la magnitud del área y de los órganos expuestos. Algunos de estos efectos, cuando la dosis se recibe en una sola exposición, se describen a continuación.

\subsection{Partes del cuerpo expuestas}

\begin{center}
\begin{tabular}{lll}
\textbf{Parte del cuerpo expuesto} & \textbf{Dosis} & \textbf{Daño} \\
Gónadas       & 50 rads   & Esterilidad temporal \\
Gónadas       & 800 rads  & Esterilidad definitiva \\
Cuero cabelludo & 500 rads  & Caída temporal del cabello \\
Cuero cabelludo & 2500 rads & Caída definitiva del cabello \\
\end{tabular}
\end{center}

\section{Efectos estocásticos y protección radiológica}

Los efectos estocásticos son aquellos que aparentemente no tienen una dosis umbral a partir de la cual se manifiestan. La probabilidad de que ocurra el efecto, independientemente de su severidad, se considera en función de la dosis; además, debe tenerse en cuenta que la mayoría de estos efectos existen normalmente en poblaciones no irradiadas. Estos efectos se pueden dividir en cuatro grupos:
\begin{enumerate}
  \item inducción de leucemia;
  \item inducción de otros cánceres;
  \item enfermedades hereditarias;
  \item acortamiento de la vida.
\end{enumerate}

El propósito de la protección radiológica es evitar los efectos determinísticos fijando límites de dosis por debajo de sus umbrales.

\section{Formas de exposición a la radiación}

Dentro de las formas de exposición a la radiación se encuentra la irradiación externa, que recibe el organismo desde fuentes exteriores al mismo. Los neutrones se consideran el tipo de radiación más peligroso, tanto por ser los más difíciles de detectar como por su efecto sobre los tejidos, ya que al interaccionar con los átomos de hidrógeno producen protones cuyo elevado poder de ionización causa un daño muy localizado.

La irradiación interna se presenta cuando la fuente se encuentra dentro del organismo. Los materiales radiactivos pueden ingresar al cuerpo por ingestión, por inhalación, por absorción a través de la piel o bien a la sangre por algún corte o herida. El daño dependerá del tipo de radiación emitida, su energía, la vida media del radionúclido, su permanencia en el organismo y el órgano en que se localice. En este caso, los emisores alfa y beta son los más peligrosos.

A continuación se anexa un cuadro que resume los efectos de la irradiación al organismo.

\begin{figure}[ht]
\centering
\includegraphics[width=\textwidth]{tablaEfectos.png}
\caption{Resumen de efectos clínicos a dosis agudas de radiación ionizante.}
\label{fig:tabla-efectos}
\end{figure}
